\documentclass[aspectratio=169,11pt]{beamer}

\usetheme{Madrid}
\usecolortheme{default}
\usefonttheme{professionalfonts}
\setbeamertemplate{navigation symbols}{}
\setbeamertemplate{footline}[frame number]

\usepackage{amsmath, amssymb, booktabs}

\title{Automation, AI, And Labor Demand}
\subtitle{Displacement, Augmentation, Scale, And New Tasks}
\author{Special Topic 5: Technology, Innovation, and Labor -- Week 2}
\date{}

\begin{document}

\begin{frame}
  \titlepage
\end{frame}

\begin{frame}{Central question}
\begin{itemize}
    \item Which tasks and workers are displaced when automation or AI changes production?
    \item Which workers are augmented because technology raises productivity in their remaining tasks?
    \item When do productivity and output-demand effects offset displacement?
    \item When does new-task creation move labor demand into different work?
\end{itemize}
\end{frame}

\begin{frame}{Labor-demand object}
\begin{itemize}
    \item The object is not technology in the abstract.
    \item The object is changing demand for tasks, skills, occupations, establishments, firms, industries, and local labor markets.
    \item A good estimate names the technology margin, the exposure mapping, the labor outcome, and the adjustment horizon.
    \item Week 2 discipline: treat measurement as part of the theory.
\end{itemize}
\end{frame}

\begin{frame}{Four demand channels}
\small
\begin{tabular}{p{0.22\linewidth} p{0.36\linewidth} p{0.26\linewidth}}
\toprule
Channel & Task effect & Labor-demand question \\
\midrule
Displacement & labor tasks move to capital/software & whose tasks shrink? \\
Augmentation & workers become more productive & whose marginal product rises? \\
Output/scale & productivity expands sales/output & does hiring recover elsewhere? \\
New tasks & new activities or occupations appear & where is labor reinstated? \\
\bottomrule
\end{tabular}
\vspace{0.35cm}
\[
\Delta L =
\Delta L^{disp} + \Delta L^{aug} + \Delta L^{scale} + \Delta L^{new}
\]
\end{frame}

\begin{frame}{Exposure is a mapping}
\[
Exposure_o = \sum_{k \in K} s_{ok}\theta_k
\]
\begin{itemize}
    \item $s_{ok}$ is the task share of task $k$ in occupation $o$.
    \item $\theta_k$ is the technology's effect on the cost or productivity of task $k$.
    \item Exposure predicts susceptibility; it is not adoption and not an effect.
    \item The outcome and counterfactual determine the labor-demand interpretation.
\end{itemize}
\end{frame}

\begin{frame}{Unit of incidence}
\small
\begin{tabular}{p{0.22\linewidth} p{0.50\linewidth} p{0.18\linewidth}}
\toprule
Unit & What the design sees & Main caution \\
\midrule
Occupation & task bundles and exposure & not adoption \\
Firm/establishment & adoption, reorganization, hiring & adopter selection \\
Industry & sector productivity and input mix & masks workers \\
Local market & equilibrium incidence across places & bundled shocks \\
\bottomrule
\end{tabular}
\vspace{0.35cm}
\begin{itemize}
    \item A technology can displace a task, raise employment at adopters, and lower wages in exposed places.
\end{itemize}
\end{frame}

\begin{frame}{Why wages and employment diverge}
\begin{itemize}
    \item Direct task displacement can lower employment in exposed work.
    \item Surviving jobs can become more skill-intensive, raising average wages.
    \item Productivity can lower prices and expand output demand.
    \item Adoption can shift rents between workers, firms, and consumers.
    \item Local labor markets adjust through mobility, entry, exit, vacancies, and occupational change.
\end{itemize}
\end{frame}

\begin{frame}{Measurement menu}
\small
\begin{tabular}{p{0.28\linewidth} p{0.34\linewidth} p{0.24\linewidth}}
\toprule
Measure & Closest object & Main warning \\
\midrule
Robot data & embodied automation adoption & narrow technology class \\
Task exposure & susceptibility of job content & not realized use \\
AI exposure & capability/task match & adoption frictions matter \\
Patents/text & invention direction & far from implementation \\
Firm adoption & realized organizational use & strong selection \\
Vacancy text & desired skill demand & not employment \\
Worker-use data & actual tool use/productivity & narrow horizon \\
\bottomrule
\end{tabular}
\end{frame}

\begin{frame}{Historical comparison}
\small
\begin{tabular}{p{0.22\linewidth} p{0.32\linewidth} p{0.32\linewidth}}
\toprule
Episode & Labor-demand lesson & AI comparison \\
\midrule
Mechanization/steam & displacement and productivity coexist & direct substitution still needs scale accounting \\
Electrification & gains require organizational redesign & AI adoption is workflow redesign, not just tool access \\
Computerization/ICT & routine tasks fall; abstract tasks rise & AI shifts the cognitive boundary of automation \\
Robotics & embodied automation gives local incidence estimates & AI data are faster but often less concrete \\
Occupational churn & new work appears slowly and unevenly & short-run AI effects miss new-task creation \\
\bottomrule
\end{tabular}
\end{frame}

\begin{frame}{What may be distinctive about AI}
\begin{itemize}
    \item It reaches some cognitive, language, coding, prediction, and coordination tasks.
    \item Diffusion can be fast where deployment costs and governance barriers are low.
    \item The treatment often works through reorganization rather than one-for-one replacement.
    \item Measurement can use text, vacancies, prompts, usage logs, and worker productivity data.
    \item General-equilibrium effects are uncertain because many sectors may adjust at once.
\end{itemize}
\end{frame}

\begin{frame}{What is not new}
\begin{itemize}
    \item Displacement anxiety has appeared in earlier technology waves.
    \item Complementary investments and adoption lags are recurring features.
    \item Skill demand changes through tasks and organization.
    \item New tasks and occupations can offset displacement, but not instantly or evenly.
    \item Incidence depends on workers, firms, places, institutions, and adjustment frictions.
\end{itemize}
\end{frame}

\begin{frame}{Research design checklist}
\begin{enumerate}
    \item Name the technology object: robot, software, AI capability, AI use, vacancy demand, or patent text.
    \item State the exposure mapping: task, occupation, firm, industry, or place.
    \item Choose the outcome: employment, wages, vacancies, sales, mobility, task mix, or job quality.
    \item State the horizon: immediate use, adoption, local adjustment, or long-run churn.
    \item Identify the equilibrium channel left outside the estimate.
\end{enumerate}
\end{frame}

\begin{frame}{Week 2 lab}
\begin{itemize}
    \item Reproduce: construct a synthetic commuting-zone robot exposure measure.
    \item Diagnose: compare employment and wage incidence across exposed local markets.
    \item Transfer: compare local robot exposure with firm AI adoption and vacancy AI demand.
    \item Goal: learn why local exposure, firm adoption, and worker use are different empirical objects.
\end{itemize}
\end{frame}

\begin{frame}{Bridge to Week 3}
\begin{itemize}
    \item Week 2 asks where labor demand moves when technology changes tasks.
    \item Week 3 asks how workers adjust once demand has moved.
    \item The next objects are skills, training, mobility, career transitions, bargaining, and welfare costs.
\end{itemize}
\end{frame}

\begin{frame}{Takeaways}
\begin{itemize}
    \item Automation and AI operate through displacement, augmentation, scale, and new tasks.
    \item Employment and wages diverge because incidence is multi-channel and multi-level.
    \item Historical comparison places AI inside a longer labor-demand literature.
    \item Measurement choices identify different margins; they are not technical afterthoughts.
\end{itemize}
\end{frame}

\end{document}
